% Einheit 1 von 4 – Zuordnungen darstellen (bank/zuordnungen/e1.jsonl, zusammenbau v0.1)
%% TODO Zweigzeile Teil 1: was hier gelernt wird (Satz in Schülersprache aus dem Einheitstitel) – Einheit 1
\einheitenkopf[e1]{Einheit 1 von 4 · Zuordnungen darstellen}
\zweigzeile{ab Kl. 6 · P10}
\setcounter{aufgabe}{7}

%% TODO \verfahren{…}: Verfahrensüberschrift in Sachsprache fehlt in der Bank (2.3 b)
%% TODO Titel: Ich-kann-Satz fehlt in der Bank (Platzhalter: Kettenname) – Nr. 8
%% TODO Anweisung: ein Satz über der ersten Teilaufgabe fehlt in der Bank – Nr. 8
\begin{aufgabe}{Darstellen}
\swa{Ein Radfahrer fährt gleichmäßig: nach 1 Stunde 15 km, nach 2 Stunden 30 km, nach 3 Stunden 45 km. Trage die Wertepaare in die Tabelle ein.}{\wertetabelleleer{Zeit in h}{Weg in km}{3}}
\swa{Trage die Punkte aus der Tabelle ins Koordinatensystem ein.}{\wertetabelle[2,4,6,8]{Zeit in h}{Weg in km}{1,2,3,4} \begin{ksys}[xmin=0,xmax=5,ymin=0,ymax=9,xlabel=Zeit in h,ylabel=Weg in km]\end{ksys}}
\swa{Trage die Punkte aus der Tabelle ins Koordinatensystem ein.}{\wertetabelle[3,6,9]{Zeit in min}{Wasser in l}{1,2,3} \begin{ksys}[xmin=0,xmax=4,ymin=0,ymax=10,xlabel=Zeit in min,ylabel=Wasser in l]\end{ksys}}
\swa{Trage die Punkte aus der Tabelle ins Koordinatensystem ein.}{\wertetabelle[1,2,3]{Anzahl Karten}{Preis in €}{2,4,6} \begin{ksys}[xmin=0,xmax=7,ymin=0,ymax=4,xlabel=Anzahl Karten,ylabel=Preis in €]\end{ksys}}
\swa{Trage die Punkte aus der Tabelle ins Koordinatensystem ein.}{\wertetabelle[3,6,9,12]{Anzahl Bücher}{Stapelhöhe in cm}{1,2,3,4} \begin{ksys}[xmin=0,xmax=5,ymin=0,ymax=13,xlabel=Anzahl Bücher,ylabel=Stapelhöhe in cm]\end{ksys}}
\swa{Trage die Punkte aus der Tabelle ins Koordinatensystem ein.}{\wertetabelle[5,10,15]{Zeit in Tagen}{gelesene Seiten}{1,2,3} \begin{ksys}[xmin=0,xmax=4,ymin=0,ymax=16,xlabel=Zeit in Tagen,ylabel=Seiten]\end{ksys}}
\swfrage{Was bleibt gleich, was ändert sich?}
\swz{Lies ab: Wie weit ist man nach 3 Stunden? Nach wie vielen Stunden sind es 8 km?}{\begin{ksys}[xmin=0,xmax=5,ymin=0,ymax=10,ablesen,xlabel=Zeit in h,ylabel=Weg in km]\gerade{2}{0}{}\end{ksys}}
\swa{Trage die Punkte aus der Tabelle ins Koordinatensystem ein. Senkrecht steht ein Kästchen für 10 l.}{\wertetabelle[20,40,60,80]{Zeit in min}{Wasser in l}{1,2,3,4} \begin{ksys}[xmin=0,xmax=5,ymin=0,ymax=100,ystep=10,xlabel=Zeit in min,ylabel=Wasser in l]\end{ksys}}
%% TODO Darstellung für schwach fehlt in der Bank (zuordnungen-e1-k1-s4-v1; Abschnitt „Für schwache Schüler“)
\swz[3]{Die Punkte gehören zur Zuordnung Anzahl Kinokarten → Preis. Dürfen sie zu einer Linie verbunden werden? Begründe.}{}
\begin{teile}
\teil Ein Eimer wird mit einem gleichmäßigen Wasserstrahl gefüllt. Welcher Graph passt zu Zeit → Füllhöhe? \\ \kreuz{steigende Gerade durch den Ursprung} \\ \kreuz{fallende Gerade} \\ \kreuz{waagerechte Linie}
\end{teile}
\swz[3]{Mia hat die Tabelle ins Koordinatensystem übertragen (Bild). Finde den Fehler.}{\wertetabelle[3,6]{Zeit in h}{Weg in km}{1,2} \begin{ksys}[xmin=0,xmax=7,ymin=0,ymax=7,xlabel=Zeit in h,ylabel=Weg in km]\punkt{3}{1}{A}\punkt{6}{2}{B}\end{ksys}}
\swa{Ein Sportverein zählt die Abonnenten seines Kanals. Trage die Werte aus der Tabelle ins Koordinatensystem ein. Die senkrechte Achse ist in Schritten von 50 beschriftet, ein Kästchen steht für 10. \hfill (P10 2025 OS)}{\wertetabelle[70,95,135,185,260]{Zeit in Monaten}{Abonnenten}{0,1,2,3,4} \begin{ksys}[xmin=0,xmax=5,ymin=0,ymax=300,ystep=50,xlabel=Zeit in Monaten,ylabel=Abonnenten]\end{ksys}}
\end{aufgabe}

%% TODO \verfahren{…}: Verfahrensüberschrift in Sachsprache fehlt in der Bank (2.3 b)
%% TODO Titel: Ich-kann-Satz fehlt in der Bank (Platzhalter: Kettenname) – Nr. 9
%% TODO Anweisung: ein Satz über der ersten Teilaufgabe fehlt in der Bank – Nr. 9
\begin{aufgabe}{Achseneinteilung}
\begin{teile}
\teil Auf einer Achse stehen die Zahlen 0, 20 und 40, je 4 Kästchen auseinander. Wie viel ist ein Kästchen wert? 1 Kästchen = \leerfeld
\end{teile}
\begin{teile}
\teil Die Werte gehen bis 60 km, die Achse hat 15 Kästchen. Welche Einteilung passt? \\ \kreuz{1 Kästchen = 1 km} \\ \kreuz{1 Kästchen = 5 km}
\end{teile}
\begin{teile}
\teil Die Werte gehen bis 40 €, die Achse hat 10 Kästchen. Welche Einteilung passt? \\ \kreuz{1 Kästchen = 2 €} \\ \kreuz{1 Kästchen = 5 €}
\end{teile}
\begin{teile}
\teil Die Werte gehen bis 9 l, die Achse hat 20 Kästchen. Welche Einteilung passt? \\ \kreuz{1 Kästchen = 1 l} \\ \kreuz{1 Kästchen = 10 l}
\end{teile}
\begin{teile}
\teil Die Werte gehen bis 300 g, die Achse hat 16 Kästchen. Welche Einteilung passt? \\ \kreuz{1 Kästchen = 10 g} \\ \kreuz{1 Kästchen = 20 g}
\end{teile}
\begin{teile}
\teil Die Werte gehen bis 70 min, die Achse hat 20 Kästchen. Welche Einteilung passt? \\ \kreuz{1 Kästchen = 2 min} \\ \kreuz{1 Kästchen = 5 min}
\end{teile}
\swfrage{Was bleibt gleich, was ändert sich?}
\swa{Wähle für die waagerechte Achse eine Einteilung und zeichne die beschriftete Achse.}{\wertetabelle[2,4,6,8,10,12]{Anzahl Hefte}{Preis in €}{1,2,3,4,5,6}}
%% TODO Darstellung für schwach fehlt in der Bank (zuordnungen-e1-k2-s3-v1; Abschnitt „Für schwache Schüler“)
\swz{Der größte Wert ist 90 km, ein Kästchen soll 5 km sein. Wie viele Kästchen braucht man? Passt es auf eine Achse mit 16 Kästchen?}{}
%% TODO Darstellung für schwach fehlt in der Bank (zuordnungen-e1-k2-s4-v1; Abschnitt „Für schwache Schüler“)
\swz{Die Werte gehen bis 800 m, die Achse hat 20 Kästchen. Wie viele Meter muss ein Kästchen mindestens tragen?}{}
\swa{Zeichne die beiden Achsen für diese Tabelle und beschrifte sie mit Größe und Einheit.}{\wertetabelle[400,800,1\,200]{Zeit in min}{Strecke in m}{5,10,15}}
\swa{Ein Regenfass enthält anfangs 90 l Wasser und läuft gleichmäßig aus. Ergänze die Tabelle. Zeichne dann auf Karopapier ein Koordinatensystem mit 20 Kästchen nach rechts und 20 Kästchen nach oben: Wähle eine Einteilung, beschrifte beide Achsen und trage die Punkte ein. \hfill (P10 2021 OS)}{\wertetabelle[,80,70,60,50,40,]{Zeit in min}{Inhalt in l}{0,5,10,15,20,25,30}}
\end{aufgabe}

%% TODO Titel: Ich-kann-Satz fehlt in der Bank (Platzhalter: Kettenname) – Nr. 10
%% TODO Anweisung: ein Satz über der ersten Teilaufgabe fehlt in der Bank – Nr. 10
\begin{aufgabe}{Formel aus Tabelle (y = 4 · x)}
\swz{Die Tabelle gehört zu einer proportionalen Zuordnung (x: Anzahl Hefte, y: Preis in €). Wie lautet die Formel?}{\wertetabelle[3,6,9]{x}{y}{1,2,3}}
\end{aufgabe}

%% TODO Titel: Ich-kann-Satz fehlt in der Bank (Platzhalter: Kettenname) – Nr. 11
%% TODO Anweisung: ein Satz über der ersten Teilaufgabe fehlt in der Bank – Nr. 11
\begin{aufgabe}{Darstellen – Fehler finden, Begründen, Darstellungswechsel, Anwendung}
\swz[3]{Lea liest am Graphen ab, wie weit sie nach 2 Stunden ist. Der Punkt liegt 5 Kästchen hoch, senkrecht steht ein Kästchen für 2 km. Lea schreibt: 5 km. Finde den Fehler.}{}
\swfrage{Kann in einer Tabelle Uhrzeit → Temperatur dieselbe Uhrzeit zweimal mit verschiedenen Temperaturen stehen? Begründe.}
\swa{Lies am Graphen die Werte für 1, 2 und 3 Stunden ab und schreibe sie als Tabelle.}{\begin{ksys}[xmin=0,xmax=4,ymin=0,ymax=6,ablesen,xlabel=Zeit in h,ylabel=Regen in mm]\gerade{1.5}{0}{}\end{ksys}}
\swz[3]{Das Diagramm zeigt die Temperatur an einem Sommertag. Lies ab: Um wie viel Uhr war es am wärmsten? Um wie viel Grad stieg die Temperatur von 8 Uhr bis dahin?}{\liniendia[ymax=25,ystep=5,yfein=1,ylabel=Temperatur in $^\circ$C]{8 Uhr/12,10 Uhr/16,12 Uhr/20,14 Uhr/23,16 Uhr/21,18 Uhr/17}}
\end{aufgabe}

\uebersichtskasten{Zuordnung: Zu jedem Wert links gehört genau ein Wert rechts. Tabelle: zwei Spalten. Graph: linker Wert waagerecht, rechter Wert senkrecht, Punkt (x | y). \\ Tabelle 2 → 5, 3 → 7,5 \quad Punkte (2 | 5), (3 | 7,5)}
